%% RevTeX 4-2 working format --- intended for a Springer economics proceedings volume
%% The aps,pra options select RevTeX typography only; they do not identify the target venue.
\documentclass[aps,pra,reprint,superscriptaddress,amsmath,amssymb,floatfix]{revtex4-2}

\usepackage{graphicx}
\usepackage{xcolor}
\usepackage{amsthm}
\usepackage[expansion=false]{microtype}
\usepackage{enumitem}
\usepackage{hyperref}
\usepackage{orcidlink}

\hypersetup{
    colorlinks=true,
    linkcolor=blue,
    filecolor=magenta,
    urlcolor=blue,
    citecolor=blue,
    pdftitle={No Trading Strategy Can Win on Every Price Path: Computability, Randomness, and the Limits of Universal Trading},
    pdfauthor={Karl Svozil},
    pdfsubject={Computational, probabilistic, and financial limits of universal trading},
    pdfkeywords={trading, no-arbitrage, diagonalization, Halting Problem, rule inference, Martin-Lof randomness, market impact}
}

%% Mathematical Environments
\theoremstyle{plain}
\newtheorem{theorem}{Theorem}
\newtheorem{lemma}[theorem]{Lemma}
\newtheorem{corollary}[theorem]{Corollary}
\newtheorem{proposition}[theorem]{Proposition}
\theoremstyle{definition}
\newtheorem{definition}{Definition}
\newtheorem{criterion}{Criterion}
\theoremstyle{remark}
\newtheorem{remark}{Remark}

\begin{document}

\title{No Trading Strategy Can Win on Every Price Path:\\
Computability, Randomness, and the Limits of Universal Trading}

\author{Karl Svozil\,\orcidlink{0000-0001-6554-2802}}
\affiliation{Institute for Theoretical Physics, TU Wien, Wiedner Hauptstrasse 8-10/136, A-1040 Vienna, Austria}
\email{karl.svozil@tuwien.ac.at}
\homepage{http://tph.tuwien.ac.at/~svozil}

\date{July 27, 2026}

\begin{abstract}
Can one algorithm make money whatever the market does? The recurring
configuration is a trader on one side and a market on the other. A responsive
computable market can use the trader's own output to build a diagonal
counterpath. If market rules can emulate universal computation, deciding
their future price events would solve the Halting Problem. At the opposite
pole, a passive Martin-L\"of-random record defeats every effective fair-game
betting test. A large or public trader may also change prices through impact
and strategic response. Repeatable success therefore requires some
exploitable restriction---trend, information, a risk premium, or another
structural insight. No-arbitrage, No-Free-Lunch averaging, and time-reversal
stress tests sharpen these limits; conditional and relative edges remain
possible.
\end{abstract}

\maketitle

%%------------------------------------------------------------------------------
\section{Introduction}
\label{sec:intro}
%%------------------------------------------------------------------------------

Throughout this paper the elementary configuration has two sides: a
\emph{trader} and a \emph{market}. At each date the trader converts the
observed price history into a position; the market supplies the next price,
and their interaction determines the gain. Depending on the model, ``market''
may mean a realized price path, a probability law over paths, or a responsive
rule that observes the trader's action. This organizing picture is not a claim
that every financial model is literally a two-person zero-sum game. It keeps
visible which powers and restrictions are assigned to each side.

The universal claim examined here has the quantifier order
\begin{equation}
  \exists\,\sigma\ \forall\,\mathcal E:
  \qquad \sigma\ \hbox{wins against the market }\mathcal E .
  \label{eq:universalquantifiers}
\end{equation}
The computational no-go result reverses it:
\begin{equation}
  \forall\,\sigma\ \exists\,\mathcal E_\sigma:
  \qquad \sigma\ \hbox{does not win against }\mathcal E_\sigma .
  \label{eq:counterquantifiers}
\end{equation}
Here the trader \(\sigma\) is a total deterministic program. After it announces
long, short, or flat, a responsive computable market can move down, up, or
nowhere, respectively. The countermarket is made from the trader's own output
by reversing that output's intended effect: this is the diagonal step.
A trader that never takes a position---that is, does not trade at all---earns zero rather than a strict loss.
The precise conclusion is therefore failure to guarantee a strictly positive gain.

Universal computation creates a second, related obstruction on the market
side. If admissible market generators can simulate arbitrary Turing machines,
a future price event can encode whether a computation halts. A program that
could infer every such rule well enough to decide every encoded event would
therefore solve the Halting Problem. This is the concrete rule-inference
barrier developed in Sec.~\ref{sec:computability}. Gold's
identification-in-the-limit framework and later work on unsolvable inductive
inference place the same boundary in a broader learning-theoretic
setting~\cite{go-67,ad-91,angluin:83,2016-AlexanderKarlSvozil-ml,moore}.
The diagonal countermarket and the Halting reduction share a self-referential
ancestry, but they prove different statements and have different quantifiers.

At the opposite pole, the market need not observe the trader at all. A passive
binary price record may be Martin-L\"of random relative to a specified
computable fair measure. Operationally, it then passes every effective null
test and permits no lower-semicomputable nonnegative fair-game capital process
to become unbounded. This does not mean that no individual move can be guessed,
or that finite lucky profits are impossible. It means that an effective trader
cannot extract a sustained fair-game advantage from that path. Algorithmic
randomness and deterministic chaos are not synonyms.

Between a passive random market and a perfectly responsive countermarket lies
the economically important case in which the trader changes the market. A
large, crowded, or publicly known strategy consumes liquidity, moves quotes,
reveals information, and attracts imitation or opposition. The price record
generated after deployment is then partly a function of the strategy itself.
This resembles an invasive classical measurement: the intervention changes
the process being measured. It is an endogeneity mechanism, not a physical
measurement postulate and not a theorem that every large trader must lose.

The practical conclusion is therefore conditional rather than nihilistic.
Repeatable success requires an \emph{inductive bias}: a restriction or
asymmetry that the trader understands and the unrestricted market class does
not assume away. A trend follower posits persistence, a contrarian posits mean
reversion, an option seller relies on a risk premium, and an informed trader
relies on information not yet fully reflected in prices. Such insights may
work in a particular environment. What cannot be supplied is an insight-free
algorithm that wins over every admissible environment.

Time reversal is the simplest transparent challenge to one such claim: hold
the trader fixed, reverse a finite market record, and run the trader causally
from the new beginning. It is one member of a countable family of computably
specified stress transformations. This countability echoes the effective-test
viewpoint of algorithmic randomness, but a reversal stress test is not itself
a Martin-L\"of test, because it carries no effective probability bound.

Here a strategy is called \emph{universal} only if, after discounting by a
specified numeraire, it produces a strictly positive terminal gain from zero
initial capital on every admissible market trajectory. The zero-cost,
self-financing, and discounted-gain qualifications are essential: a positive
initial investment in a bank account has positive terminal wealth, but it is
not an arbitrage and does not outperform its financing benchmark.

The mathematical study of unpredictable prices begins with Bachelier's 1900
thesis~\cite{Bachelier1900}, continues through Samuelson's martingale
formulation~\cite{Samuelson1965}, and informs Fama's Efficient Market
Hypothesis (EMH)~\cite{Fama1970}. The EMH is an economic claim about the
incorporation of information into prices; it is not, by itself, a theorem that
no individual strategy can win unconditionally.

Section~\ref{sec:computability} gives the diagonal counterpath, two explicit
Halting reductions, and the Martin-L\"of, computable, and Schnorr randomness
hierarchy. Section~\ref{sec:criterion} develops time reversal and other
computable stress transformations. Sections~\ref{sec:noarb} and
\ref{sec:nfl} give the separate no-arbitrage and No-Free-Lunch boundaries.
The remaining sections delimit the social-choice analogy, analyze the Wheel
Options Strategy, distinguish Cover's benchmark-relative universality, and
state the honest conditional scope of strategies that work \emph{for all
practical purposes} (FAPP).

%%------------------------------------------------------------------------------
\section{The Trader and the Market: Computability and Randomness}
\label{sec:computability}
%%------------------------------------------------------------------------------

We now formalize the trader--market pair. The trader is a program
\(M_\sigma\) mapping each observed finite price history to a position; the
market is the environment that supplies the next price. Three specifications
of the market side lead to three different boundaries: a responsive computable
market yields a tailored counterpath, a computationally universal class of
market generators makes unrestricted future-event inference undecidable, and
a passive algorithmically random record defeats effective fair betting.

\subsection{A Responsive Market: The Diagonal Counterpath}

First discard the assumption of an exogenous price process and let the market
react after observing the trader's current position. This is a worst-case,
perfect-information model. Let \(M_\sigma\) be a \emph{total} computable map
from every finite computable price history \(H_t\) to a rational position
\(w_t\). Totality models the practical requirement that a deployed algorithm
return an action before a deadline.

\begin{theorem}[Computable Adaptive Counterpath]
  \label{thm:counterpath}
  For every total deterministic computable strategy \(M_\sigma\) and every
  finite horizon \(T\), there exists a positive computable price path
  \(P_{\mathrm{adv}}\) on which its self-financing gain
  \[
    G_T=\sum_{t=0}^{T-1}w_t(P_{t+1}-P_t)
  \]
  is non-positive. It is strictly negative if the strategy takes a nonzero
  position at least once.
\end{theorem}

\begin{proof}
Choose computable \(P_0>0\) and rational \(0<\epsilon<1\). At time \(t\),
compute \(w_t=M_\sigma(H_t)\) and set
\begin{equation}
P_{t+1}=
\begin{cases}
P_t(1-\epsilon),&w_t>0,\\
P_t(1+\epsilon),&w_t<0,\\
P_t,&w_t=0.
\end{cases}
\label{eq:adversary}
\end{equation}
The resulting prices remain positive and computable. In each period,
\(w_t(P_{t+1}-P_t)\leq0\), with strict inequality when \(w_t\neq0\).
Summing proves the claim.
\end{proof}

The construction is diagonal in the following precise operational sense: the
market takes the strategy's own current output and applies a fixed-point-free
sign reversal to it. Long is answered by down, short by up, and neutrality by
zero. The candidate rule is therefore defeated on an input made from its own
answer, just as a diagonal argument defeats a purported universal rule by
turning that rule against itself. No enumeration of all machines is required.
Its quantifiers are
\(\forall M_\sigma\,\exists P_{\mathrm{adv}}\): the counterpath may depend on
the strategy. There need not be one finite path on which every strategy loses;
for example, a permanently long and a permanently short position cannot both
lose on the same one-step price move.

\subsection{Universal Computation and the Halting Barrier}

\begin{theorem}[Undecidability of Eventual Profit]
  \label{thm:profit-undecidable}
  Fix the positive computable price path \(P_t=t+1\), \(t\in\mathbb{N}\).
  There is no algorithm that, given the code of an arbitrary total computable
  trading program \(M_\sigma\), always decides whether
  \[
    \exists T\geq1:\qquad
    G_T(M_\sigma;P)=
    \sum_{t=0}^{T-1}w_t(P_{t+1}-P_t)>0.
  \]
\end{theorem}

\begin{proof}
Assume such a profit certifier exists. Given an arbitrary Turing machine
\(A\) and input \(x\), construct a trading program \(M_{A,x}\) as follows.
At date \(t\), simulate \(A(x)\) for \(t+1\) steps. If it has halted within
those steps, set \(w_t=1\); otherwise set \(w_t=0\). Every date requires only
a finite simulation, so \(M_{A,x}\) is total. Along \(P_t=t+1\), every long
position earns one unit in the next period. Hence \(G_T>0\) for some \(T\) if
and only if \(A(x)\) eventually halts. A certifier for eventual trading profit
would therefore decide the Halting Problem, which is impossible
\cite{turing-36}.
\end{proof}

The reduction can also be placed directly on the market side, which is the
form most closely connected to general rule inference.

\begin{theorem}[Undecidability of a Market Event]
  \label{thm:market-undecidable}
  There is no algorithm that, given the code of an arbitrary total computable
  positive-price generator \(G\), always decides whether its generated path
  \(P^G\) ever satisfies \(P_t^G>P_0^G\).
\end{theorem}

\begin{proof}
Given a Turing machine \(A\) and input \(x\), define the total generator by
\(P_0^{A,x}=1\) and, for \(t\geq1\),
\begin{equation}
  P_t^{A,x}
  =1+\mathbf{1}\{A(x)\text{ halts within }t\text{ steps}\},
  \qquad t\geq1 .
  \label{eq:market-halting}
\end{equation}
Each quoted price requires only a finite \(t\)-step simulation, so the
generator is total and computable. Its initial price is one, and it ever quotes
a higher price if and only if \(A(x)\) halts. A universal decider for the stated
market event would therefore decide the Halting Problem.
\end{proof}

Theorem~\ref{thm:market-undecidable} makes the assumption of universal
computation precise: the admissible language of market generators must be rich
enough to emulate an arbitrary machine and expose the result through an
observable price event. Under that condition, no computable rule-inference
procedure can uniformly settle all future events of all such markets. This is
a concrete market version of the limits studied in identification in the limit
and in unsolvable inductive inference
\cite{go-67,blum75blum,angluin:83,ad-91,2016-AlexanderKarlSvozil-ml,moore}.
It does not deny that restricted parametric, stochastic, or otherwise
structured model classes can be learned.

The three computational claims are related but not identical.
Theorem~\ref{thm:counterpath} has quantifiers
\(\forall M_\sigma\,\exists P_{\mathrm{adv}}\) and constructs a countermarket
from the action of a fixed trader; that construction only needs enough
computation to run that trader. Theorem~\ref{thm:profit-undecidable} embeds an
arbitrary computation in the trader while fixing a very simple rising market.
Theorem~\ref{thm:market-undecidable} embeds it in the market and rules out a
general future-event inference procedure. For a specified finite horizon, by
contrast, total programs can simply be run and their realized gains and prices
computed.

\begin{corollary}[Maximin Value]
  \label{cor:maximin}
  In the finite-horizon game of Theorem~\ref{thm:counterpath}, if the neutral
  strategy is allowed, then
  \[
    \sup_{M_\sigma}\inf_{P_{\mathrm{adv}}}G_T(M_\sigma,P_{\mathrm{adv}})=0.
  \]
\end{corollary}

\begin{proof}
The neutral strategy guarantees zero, while
Theorem~\ref{thm:counterpath} gives every strategy a counterpath with gain at
most zero.
\end{proof}

\begin{proposition}[Private Randomization Does Not Create a Guarantee]
  \label{prop:randomized}
  Suppose a randomized trader's current coin toss is hidden from the
  environment. There still exists a passive computable distribution of next
  returns under which every integrable strategy has conditional expected gain
  zero.
\end{proposition}

\begin{proof}
Let the next price increment be \(+\epsilon P_t\) or \(-\epsilon P_t\) with
equal probability, independently of the trader's private draw. Conditional on
the current information and position,
\(\mathbb{E}[w_t(P_{t+1}-P_t)\mid\mathcal{F}_t]=0\).
Summing gives expected gain zero.
\end{proof}

Randomization therefore prevents a code-observing opponent from mechanically
opposing the unseen current draw, but it does not yield strictly positive
expected gain in every environment.

\subsection{The Responsive Market as a Game}

Theorem~\ref{thm:counterpath} can be read as a finite perfect-information
zero-sum game: the trader announces \(w_t\), the environment chooses the next
return, and the terminal payoff is \(G_T\). Its value is zero by
Corollary~\ref{cor:maximin}; backward induction, not a diagonal enumeration,
identifies the environment's response. For infinite perfect-information games,
Borel determinacy guarantees that one player has a pure winning strategy when
the winning set is Borel~\cite{Martin1975}, but the theorem does not identify
the winning player without analysis of the payoff set. Nor does it cover a
trader's hidden random bits: that is an imperfect-information game whose
relevant conclusion here is the mixed value zero in
Proposition~\ref{prop:randomized}.

\subsection{A Passive Market: Martin-L\"of Randomness}

The market need not observe or answer the trader. At the opposite pole,
suppose an effective coding of price directions is Martin-L\"of random
relative to a specified computable fair measure. Here ``random'' means
algorithmically random, not merely sensitive dependence in deterministic
chaos. Such a sequence may still be guessed correctly on many individual
dates. The exact claim is that it passes every effective null test, or
equivalently that no lower-semicomputable nonnegative fair-game capital process
becomes unbounded on it.

In the fair binary market, a nonnegative capital process
\(K\colon\{0,1\}^{<\mathbb{N}}\to\mathbb{R}_{\geq0}\) is a martingale when
\begin{equation}
  K(s)=\frac{K(s0)+K(s1)}2.
  \label{eq:computablemartingale}
\end{equation}

\begin{theorem}[Ville's Inequality]
  \label{thm:ville}
  If \(K_0=1\) is a nonnegative martingale under the fair-coin measure, then
  for every \(\lambda>0\),
  \begin{equation}
    \mathbb{P}\!\left(\sup_{n\geq0}K_n\geq\lambda\right)
    \leq\frac1\lambda.
    \label{eq:ville}
  \end{equation}
\end{theorem}

Ville's inequality follows by stopping the nonnegative capital process when it
first crosses \(\lambda\) and applying optional stopping, followed by a
limit~\cite{Ville1939}. In particular, the set on which a fixed nonnegative
martingale becomes unbounded has probability zero.

To state the effective refinement correctly, let \(\lambda\) denote fair-coin
measure on Cantor space \(\{0,1\}^{\mathbb N}\). A Martin-L\"of test is a
uniformly computably enumerable sequence of open sets
\((U_k)_{k\geq1}\) satisfying \(\lambda(U_k)\leq2^{-k}\). A sequence is
\emph{Martin-L\"of random} if it avoids \(\bigcap_k U_k\) for every such
test~\cite{MartinLof1966}. A sequence is \emph{computably random} if no total
computable nonnegative martingale succeeds on it by becoming unbounded. A
\emph{Schnorr-random} sequence passes every Martin-L\"of test whose component
measures are uniformly computable. These are distinct notions, with strict
implications
\begin{equation}
  \begin{aligned}
    \text{Martin-L\"of random}
      &\ \Longrightarrow\ \text{computably random}\\
      &\ \Longrightarrow\ \text{Schnorr random}.
  \end{aligned}
  \label{eq:randomhierarchy}
\end{equation}
Neither converse holds~\cite{schnorr1,Nies2009,DH}.

\begin{theorem}[Effective Martingale Barrier]
  \label{thm:effective}
  A binary sequence \(x\) is Martin-L\"of random if and only if no
  lower-semicomputable nonnegative supermartingale succeeds on \(x\) by
  unbounded capital. Consequently, the success set of each such betting
  strategy is effectively \(\lambda\)-null, and the set of paths on which all
  of them remain bounded has fair-coin measure one.
\end{theorem}

\begin{proof}
Normalize a lower-semicomputable nonnegative supermartingale \(d\) by
\(d(\varnothing)\leq1\), and let
\[
  U_k=\{x:\exists n,\ d(x{\upharpoonright}n)\geq2^k\}.
\]
Lower semicomputability makes \(U_k\) computably enumerable and open, while
Ville's inequality gives \(\lambda(U_k)\leq2^{-k}\). Hence
\((U_k)\) is a Martin-L\"of test and no Martin-L\"of-random sequence permits
\(d\) to become unbounded. Conversely, every Martin-L\"of test can be converted
into a lower-semicomputable supermartingale that succeeds on its intersection;
this is the martingale characterization of Martin-L\"of
randomness~\cite{calude:02,Nies2009}.
\end{proof}

The same construction extends from \(\lambda\) to any computable probability
measure \(\mu\) with the supermartingale condition and cylinder probabilities
defined relative to \(\mu\). This is the effective analogue most directly
applicable to a computable martingale measure.

Theorem~\ref{thm:effective} contains the proposed ``typical path'' claim in its
correct form. Under a computable fair measure, a measure-one class of paths
simultaneously blocks unbounded success by every effective nonnegative
fair-game capital process. It does \emph{not} say that every such path gives
every strategy a finite loss, nor does it apply unchanged to strategies with
unrestricted borrowing or to non-martingale physical measures.
There is also a direct prediction reading. If a total computable next-bit
predictor maintained a long-run accuracy strictly above \(1/2\), then betting
a sufficiently small fixed fraction of capital on each prediction would
produce an unbounded computable martingale. Thus no such sustained directional
advantage is possible on a Martin-L\"of-random fair-coin path, although finite
streaks and isolated correct forecasts remain possible.

Theorem~\ref{thm:counterpath} is different: its market is tailored to one
deterministic trader and gives that trader a finite non-positive gain. A
Martin-L\"of-random market need not react to the trader and simultaneously
passes all effective tests. The former is a worst-case construction; the
latter is a typical-path statement relative to a specified computable
measure. Game-theoretic probability develops this pathwise betting viewpoint
systematically~\cite{ShaferVovk2001}.

\subsection{What Rice's Theorem Does and Does Not Say}

Rice's theorem states that every nontrivial extensional property of arbitrary
partial computable functions is undecidable~\cite{Rice-1953}. It therefore
limits fully general program certification; even deciding whether arbitrary
code is total is impossible in general~\cite{turing-36}. It does \emph{not}
imply that the loss-producing inputs of a fixed total finite-horizon strategy
are undecidable. On a specified finite computable path, such a program can be
simulated and its gain calculated. No probability claim about market failure
sets follows from Rice's theorem alone. Theorem~\ref{thm:profit-undecidable}
supplies the exact statement needed here by an explicit Halting-Problem
reduction: it is the unbounded question of \emph{eventual} profit, not the
evaluation of a completed finite trade record, that is undecidable in general.

%%------------------------------------------------------------------------------
\section{The Time-Reversal Stress Test}
\label{sec:criterion}
%%------------------------------------------------------------------------------

The trader--market viewpoint suggests holding the trader fixed while
transforming the market record. Among computably specified stress
transformations, time reversal is perhaps the simplest.

Let \(\mathcal{M}\) be a space of strictly positive paths
\(P\colon[0,T]\to\mathbb{R}_{>0}^n\). The time reversal of \(P\) is
\begin{equation}
  \widetilde{P}(t)=P(T-t), \qquad t\in[0,T].
  \label{eq:timereversal}
\end{equation}
For a causal strategy \(\sigma\), let \(\Pi_T[\sigma;P]\) denote its terminal
discounted gain from zero initial capital when evaluated on \(P\).

\begin{definition}[Pathwise Universal Strategy]
\label{def:universal}
  A causal strategy \(\sigma\) is \emph{pathwise universal} on
  \(\mathcal{M}\) if it is self-financing and admissible and satisfies
  \(\Pi_T[\sigma;P]>0\) for every \(P\in\mathcal{M}\).
\end{definition}

\begin{criterion}[Reversal Falsification Test]
  \label{crit:tr}
  Suppose \(\mathcal{M}\) is closed under time reversal. If there exists
  \(P\in\mathcal{M}\) such that
  \begin{equation}
    \min\{\Pi_T[\sigma;P],\Pi_T[\sigma;\widetilde P]\}\leq 0,
    \label{eq:reversaltest}
  \end{equation}
  then \(\sigma\) is not pathwise universal on \(\mathcal{M}\).
\end{criterion}

\begin{proof}
Both \(P\) and \(\widetilde P\) belong to \(\mathcal{M}\). A non-positive gain
on either member of the pair contradicts Definition~\ref{def:universal}.
\end{proof}

Criterion~\ref{crit:tr} is deliberately a falsification device, not an
independent impossibility theorem. Reversing a profitable path does not in
general negate the profit of a causal strategy, because the strategy can take
different positions on the reversed history. The value of the test is that it
generates a paired stress scenario whenever the modeled path class is
reversal-closed.

\begin{remark}[Computable Stress Tests Beyond Reversal]
\label{rem:stressfamily}
Time reversal is only the first member of a much larger battery. Depending on
the admissible market class, a computable probe can reverse return signs,
splice trend and anti-trend regimes, insert jumps or observation delays,
rescale volatility, or add financing and transaction costs. Because programs
have finite descriptions, computably specified probes form an at most
countable family, and countably infinite test batteries can be constructed.
Each transformed record must still belong to the market class against which
the strategy's claim is made.

The word ``test'' is used here in the empirical sense of stress and
falsification. A Martin-L\"of test is a different technical object: a uniformly
computably enumerable sequence of open sets with effective measure bounds.
The shared idea is to confront an algorithm with many effective challenges;
time reversal by itself neither defines a null set nor certifies randomness.
Passing it establishes limited robustness, never universality.
\end{remark}

\begin{remark}[Filtrations and Time Reversal]
In stochastic finance, the set of admissible trajectories is not automatically
closed under time reversal. The process \(P(T-t)\) is generally not adapted to
the original forward filtration, and its semimartingale property under a
reversed filtration requires additional hypotheses; time reversal of
diffusions is a separate theorem, not a formal substitution
\(t\mapsto T-t\)~\cite{HaussmannPardoux1986}.

A simple Black--Scholes calculation illustrates both the useful case and its
limit. If
\[
  X_t=\log(S_t/S_0)=mt+\sigma W_t,
\]
then the reversed increment \(Y_t=X_{T-t}-X_T\), relative to its reversed
filtration, has the law of a Brownian motion with drift \(-m\) and volatility
\(\sigma\). Thus a log-price model class containing both drift signs is closed
in law under reversed increments. Under the physical measure,
\(m=\mu-\sigma^2/2\); under a risk-neutral measure the discounted stock is a
martingale, but its logarithm still has It\^o drift \(-\sigma^2/2\).
Consequently, neither an ELMM nor a driftless arithmetic Brownian special case
establishes detailed balance for a general financial model. The stress test
remains pathwise; the rigorous probabilistic constraint is given next.
\end{remark}

%%------------------------------------------------------------------------------
\section{The No-Arbitrage Foundation}
\label{sec:noarb}
%%------------------------------------------------------------------------------

Here the market is no longer an opponent choosing a response after seeing the
trader's action. Instead, financial admissibility and no-arbitrage delimit
which trader--market pairs are coherent.

\subsection{Discounted Gains and Admissibility}

Let the market be defined on a filtered probability space
\((\Omega,\mathcal{F},\{\mathcal{F}_t\}_{t\in[0,T]},\mathbb{P})\).
Let \(B_t>0\) be a numeraire and \(S_t\) the vector of traded asset prices.
Write \(\overline S_t=S_t/B_t\) for discounted prices. A predictable,
\(\overline S\)-integrable strategy \(H\) has discounted self-financing gain
\begin{equation}
  G_t=(H\mathbin{\cdot}\overline S)_t,\qquad G_0=0.
  \label{eq:gain}
\end{equation}
It is \emph{admissible} if there is a finite constant \(a\) such that
\begin{equation}
  G_t\geq -a\quad\text{for all }t\in[0,T],\quad\mathbb{P}\text{-a.s.}
  \label{eq:admissible}
\end{equation}
This uniform lower bound rules out doubling systems financed by unbounded
credit.

\begin{theorem}[Universality Implies Classical Arbitrage]
  \label{thm:arbitrage}
  If an admissible self-financing strategy \(H\) with zero initial capital
  satisfies \(G_T(\omega)>0\) for every \(\omega\in\Omega\), then \(H\) is a
  classical arbitrage.
\end{theorem}

\begin{proof}
The pathwise assumption implies \(G_T\geq0\) almost surely and
\(\mathbb{P}(G_T>0)=1\). Together with zero initial cost, self-financing, and
admissibility, these are stronger than the usual classical-arbitrage
conditions \(\mathbb{P}(G_T\geq0)=1\) and
\(\mathbb{P}(G_T>0)>0\).
\end{proof}

\begin{remark}[Terminology and Strength]
The pathwise condition is sometimes called a sure or strong arbitrage, but it
is not synonymous with \emph{arbitrage of the first kind}. The latter is the
condition dual to No Unbounded Profit with Bounded Risk (NUPBR/NA1) and is a
distinct, weaker viability concept~\cite{Kardaras2012}. Classical
no-arbitrage alone already excludes the universal strategy in
Theorem~\ref{thm:arbitrage}; NFLVR is invoked below because it supplies the
standard martingale-measure formulation.
\end{remark}

\subsection{NA, NUPBR, and Benchmark-Relative Gains}

The relevant no-arbitrage notions do not form the simple chain sometimes
suggested in informal discussions. In the standard semimartingale setting,
NFLVR combines classical NA with NUPBR, whereas NA and NUPBR are in general
logically distinct. NUPBR is equivalent, under standard portfolio hypotheses,
to the existence of a numeraire portfolio whose positive relative-wealth
processes are supermartingales~\cite{KaratzasKardaras2007,Kardaras2012}.
It may hold even when the stronger NFLVR condition and an ELMM fail.

This distinction prevents two overstatements. First, Theorem~\ref{thm:arbitrage}
uses only classical NA; NUPBR alone is not silently substituted for it. If a
zero-cost gain were nonnegative at all intermediate times and positive at
terminal time, scaling would violate NUPBR, but terminal pathwise positivity
together with only a strategy-specific lower bound does not by itself supply
that scaling argument. Second, benchmark and stochastic-portfolio approaches
can generate gains \emph{relative} to a numeraire or market portfolio without
producing zero-cost positive gain on every path. Relative arbitrage in diverse
equity markets and real-world benchmark pricing therefore do not contradict
the theorem~\cite{FernholzKaratzasKardaras2005,PlatenHeath2010}.

\subsection{The Martingale Constraint}

For a locally bounded semimartingale price process, the
Delbaen--Schachermayer form of the First Fundamental Theorem states that No
Free Lunch with Vanishing Risk (NFLVR) is equivalent to the existence of an
equivalent local martingale measure (ELMM)
\(\mathbb{Q}\approx\mathbb{P}\) under which \(\overline S\) is a local
martingale~\cite{Harrison1981,Delbaen1994}. For general unbounded
semimartingales the corresponding formulation uses an equivalent
\(\sigma\)-martingale measure~\cite{DelbaenSchachermayer1998}; the corollary below
retains the locally bounded setting.

\begin{corollary}[No Universal Zero-Cost Strategy]
  \label{cor:noarb}
  If an ELMM \(\mathbb{Q}\) exists, no admissible self-financing strategy with
  zero initial capital can have \(G_T>0\) on every path.
\end{corollary}

\begin{proof}
Under \(\mathbb{Q}\), the admissible stochastic integral
\(G=H\mathbin{\cdot}\overline S\) is a local martingale bounded from below and
therefore a supermartingale. Hence
\(\mathbb{E}_{\mathbb{Q}}[G_T]\leq G_0=0\). If \(G_T>0\) on every path, then
equivalence gives \(G_T>0\), \(\mathbb{Q}\)-a.s., and consequently
\(\mathbb{E}_{\mathbb{Q}}[G_T]>0\), a contradiction.
\end{proof}

The ELMM conclusion concerns discounted gains under the pricing measure
\(\mathbb{Q}\). It does not say that every risky strategy has non-positive
expected return under the physical measure \(\mathbb{P}\); positive physical
expected returns may compensate risk.

\subsection{Win Rate Is Not Expected Value}

The probability of closing a trade at a profit can greatly exceed \(1/2\)
even in a fair market. Optional stopping gives the precise distinction.

\begin{proposition}[Asymmetric Barriers in a Fair Random Walk]
  \label{prop:optional}
  Let \(X_n\) be a simple symmetric random walk with \(X_0=0\), and for positive
  integers \(a,b\) let
  \[
    \tau=\inf\{n\geq0:X_n=a\text{ or }X_n=-b\}.
  \]
  Then
  \begin{equation}
    \mathbb{P}(X_\tau=a)=\frac{b}{a+b},
    \qquad
    \mathbb{E}[X_\tau]=0.
    \label{eq:gambler}
  \end{equation}
\end{proposition}

\begin{proof}
The stopped process \(X_{n\wedge\tau}\) is a bounded martingale, so Doob's
optional-stopping theorem applies~\cite{Doob1953}. If
\(p=\mathbb{P}(X_\tau=a)\), then
\(0=\mathbb{E}[X_\tau]=ap-b(1-p)\), which gives
\(p=b/(a+b)\).
\end{proof}

With a take-profit of \(a=1\) and a stop-loss of \(b=9\), the win probability
is \(0.9\), but the expected payoff is zero: nine frequent unit gains are
offset by one nine-unit loss in expectation. Thus prediction accuracy, trade
win rate, and expected discounted profit are three different quantities.

%%------------------------------------------------------------------------------
\section{The Combinatorial Perspective}
\label{sec:nfl}
%%------------------------------------------------------------------------------

Here the opposing side is an ensemble of market records rather than an
adaptive market. The optimization NFL theorem and the elementary prediction
result needed for market-direction claims should be distinguished.

\begin{theorem}[Optimization NFL, Wolpert--Macready]
  \label{thm:nfl}
  For finite search and cost spaces, when performance is averaged uniformly
  over all objective functions, any two non-revisiting black-box search
  algorithms have identical averaged
  performance~\cite{Wolpert:1997:NFL:2221336.2221408}.
\end{theorem}

The theorem concerns oracle-based optimization over function spaces. It does
not, by itself, imply a statement about self-financing gains on stochastic
price paths. A sharpening makes its structural assumption explicit: for a
uniform distribution on a finite function class, the NFL equality for all
non-revisiting algorithms holds precisely when the class is closed under
permutations of the search domain~\cite{SchumacherVoseWhitley2001}. Financial
path classes are generally not permutation-closed: temporal order, continuity,
financing, and no-arbitrage restrictions distinguish one permutation from
another. Thus a financial edge, when present, resides in exactly such
symmetry-breaking structure and the inductive bias used to model it.

For binary directional prediction the relevant result is simpler and can be
derived directly.

\begin{proposition}[Uniform Binary Prediction]
  \label{prop:binary}
  Let \(X_1,\ldots,X_T\) be uniformly distributed on \(\{0,1\}^T\). A causal
  deterministic predictor chooses
  \(\widehat X_t=f_t(X_1,\ldots,X_{t-1})\). Its mean accuracy
  \[
    A_T=\frac1T\sum_{t=1}^T
      \mathbf{1}\{\widehat X_t=X_t\}
  \]
  satisfies \(\mathbb{E}[A_T]=1/2\). The same holds for a randomized predictor
  whose private randomization is independent of the next bit.
\end{proposition}

\begin{proof}
Conditional on every history \(X_1,\ldots,X_{t-1}\), the next bit \(X_t\) is
an independent fair bit. Therefore
\(\mathbb{P}(\widehat X_t=X_t\mid X_1,\ldots,X_{t-1})=1/2\).
Linearity of expectation gives the result. Conditioning additionally on the
predictor's private randomization proves the randomized case.
\end{proof}

Consequently, no predictor can achieve expected directional accuracy greater
than \(1/2\) under the uniform distribution over all binary sequences; in
particular, none can exceed \(1/2\) on every sequence. This is an NFL symmetry
statement, not a model of actual markets. Financial sign sequences are not
uniform in general, and a directional forecast does not determine trading
profit because payoff sizes, financing, and stopping rules matter
(Proposition~\ref{prop:optional}).

The defensible practical conclusion is conditional: every repeatably
successful predictor or strategy exploits non-uniform structure in the
physical measure \(\mathbb{P}\). This is the trader's indispensable insight or
inductive bias---a trend, mean reversion, a risk premium, information, or an
institutional constraint that breaks the NFL symmetry. Its edge is evidence
about a restricted environment, not universal algorithmic superiority.

%%------------------------------------------------------------------------------
\section{Analogy to Social Choice: Scope and Limits}
\label{sec:arrow}
%%------------------------------------------------------------------------------

Arrow's Impossibility Theorem concerns social welfare functions on at least
three alternatives. Under unrestricted preference profiles, collective
rationality, Pareto unanimity, and independence of irrelevant alternatives,
the aggregation rule must be dictatorial~\cite{Arrow1963}. Its assumptions and
objects differ from those of a trading model.

There is nevertheless a useful structural comparison:
\begin{description}
  \item[Universal domain.] Arrow quantifies over all admissible preference
  profiles; pathwise trading universality quantifies over all admissible price
  paths.
  \item[Jointly incompatible demands.] Arrow's fairness and consistency
  requirements cannot all be met by a non-dictatorial rule. In trading,
  zero-cost strict profit on every path is incompatible with no-arbitrage.
  \item[No theorem transfer.] Neither result proves the other. The parallel is
  a shared warning that unrestricted domains make apparently natural
  universal requirements inconsistent.
\end{description}

Yanofsky shows how many self-referential arguments can be organized using
fixed points and fixed-point-free transformations~\cite{Yanofsky-BSL:9051621}.
The response in Eq.~\eqref{eq:adversary} anti-aligns returns with every
nonzero position, but it admits the neutral fixed point \(w_t=0\). That fixed
point yields zero gain rather than strict profit. The relevant obstruction is
therefore the absence of a \emph{strictly profitable} fixed response, not the
absence of every fixed point.

%%------------------------------------------------------------------------------
\section{Case Study: The Wheel Options Strategy}
\label{sec:wheel}
%%------------------------------------------------------------------------------

The Wheel sells a cash-secured put and, after assignment, sells a covered call
against the acquired stock. Its economic exposure follows directly from
put--call parity. Ignoring dividends and financing for the terminal-payoff
identity,
\begin{equation}
  S_T-(S_T-K)^+
  =K-(K-S_T)^+.
  \label{eq:putcall}
\end{equation}
The left-hand side is a covered-call payoff; the right-hand side is cash
paying \(K\) plus a short put. After the corresponding time-zero costs are
matched, a covered call and a cash-secured put are equivalent European
payoffs. Rolling the Wheel therefore repeatedly maintains a put-like,
short-downside-convexity exposure rather than creating directionless income.
Equivalently,
\begin{equation}
  \Phi_K(S_T)
  =S_T-(S_T-K)^+
  =
  \begin{cases}
    S_T,&S_T<K,\\
    K,&S_T\geq K,
  \end{cases}
  \label{eq:wheelpiecewise}
\end{equation}
which makes the one-for-one downside and capped upside explicit.

Option-selling returns may nevertheless have a positive physical mean. The
difference between risk-neutral and physical expectations of volatility is a
variance risk premium; option-market evidence documents compensation for
bearing volatility and jump risk~\cite{BakshiKapadia2003,CarrWu2009}. This is
an empirical, distribution-dependent explanation for some Wheel income, not a
pathwise identity: the sign and magnitude depend on prices, strikes,
transaction costs, and the physical law.

\subsection{Failure I: Catastrophic Drop}

If the underlying crashes after put assignment, the short call liability
decreases, which benefits the call writer, but the loss on the stock dominates
and the combined covered-call position can suffer a large net loss.
Equation~\eqref{eq:putcall} makes this explicit: below \(K\), the payoff falls
one-for-one with \(S_T\). A reversed rally/crash pair can be used as the stress
test in Criterion~\ref{crit:tr}, but time reversal alone is not the proof; the
payoff identity is.

\subsection{Failure II: Persistent Decline}

Repeated option premiums offset part of a persistent decline but do not remove
the renewed short-put exposure. The loss is path-dependent rather than
generally linear: strikes, maturities, volatility, assignment dates, and
position sizing determine the sequence of gains.

\subsection{Failure III: Violent Breakout and Benchmark Underperformance}

In a large rally, shares may be called away at the call strike. The Wheel can
still realize a positive absolute profit, so opportunity cost alone is not a
counterexample to strict positive P\&L. It is instead a failure relative to a
buy-and-hold benchmark: Eq.~\eqref{eq:putcall} caps the terminal payoff at
\(K\). No-arbitrage constrains the price of this payoff under
\(\mathbb{Q}\); it does not say that its physical expected return must be zero.

\subsection{Failure IV: Absorbing Ruin Under Capital Constraints}

A fully cash-secured, unlevered cycle has bounded loss, but an underlying that
falls to zero can consume almost all committed capital net of premium.
Leverage makes the absorbing nature of ruin explicit. If wealth evolves as
\(W_{t+1}=W_t(1+fR_{t+1})\) and an admissible return satisfies
\(1+fR_{t+1}=0\), then \(W_{t+1}=0\); a self-financing strategy cannot recover
without an external capital injection. This is a wealth-dynamics fact, not a
consequence of undecidability.

\begin{table*}[ht]
\caption{\label{tab:failuremodes}Wheel failure modes and the mathematical mechanism illustrating each one.}
\begin{ruledtabular}
\begin{tabular*}{\textwidth}{@{\extracolsep{\fill}}llll}
\textbf{Failure} & \textbf{Trajectory Class} & \textbf{Relevant Result} & \textbf{Mechanism} \\
\colrule
I: Crash & Large downside move & Eq.~\eqref{eq:putcall} & Put-like downside \\
II: Bleed & Persistent negative drift & Repeated payoff exposure & Premiums do not insure the stock \\
III: Breakout & Large upside move & Eq.~\eqref{eq:putcall} & Capped benchmark return \\
IV: Ruin & Capital exhaustion & Wealth recursion & Zero wealth is absorbing \\
\end{tabular*}
\end{ruledtabular}
\end{table*}

Table~\ref{tab:failuremodes} is a payoff-based classification, not a claim
that each trajectory has positive probability under every continuous price
model. Corollary~\ref{cor:noarb} excludes strict positive discounted gain on
all paths, but probabilities of particular failure classes require a specified
physical measure.

Repeated option income also raises a time-versus-ensemble distinction. For
multiplicative wealth \(W_T=W_0\prod_{t=1}^T(1+fR_t)\), the realized
per-period growth is
\[
  \frac1T\log\frac{W_T}{W_0}
  =\frac1T\sum_{t=1}^T\log(1+fR_t),
\]
not the logarithm of an ensemble-average payoff. Frequent small gains can
coexist with poor long-run growth when rare losses are large, a point emphasized
in discussions of non-ergodic economic growth~\cite{Peters2019}. This does not
replace expected-utility analysis; it identifies which long-run observable is
being optimized.

%%------------------------------------------------------------------------------
\section{Universal Portfolios and Distance from Universality}
\label{sec:cover}
%%------------------------------------------------------------------------------

The word ``universal'' also appears in a celebrated positive theorem.
For price-relative vectors \(x_t\in\mathbb{R}_{>0}^m\), a constant-rebalanced
portfolio \(b\) has wealth
\[
  S_T(b)=\prod_{t=1}^T b^\mathsf{T}x_t.
\]
Cover's universal portfolio is nonanticipating and, under the conditions of
his theorem, has the same asymptotic exponential growth rate as the best
constant-rebalanced portfolio chosen in hindsight on every individual bounded
sequence~\cite{Cover1991}. In regret notation,
\begin{equation}
  \frac1T\left[
  \log\max_{b}S_T(b)-\log S_T^{\mathrm{UP}}
  \right]\longrightarrow0.
  \label{eq:cover}
\end{equation}

Equation~\eqref{eq:cover} is a relative guarantee, not an absolute-profit
guarantee. If every constant-rebalanced portfolio loses money on a path, the
universal portfolio may lose as well while retaining vanishing per-period
regret. Cover's result therefore does not contradict
Corollary~\ref{cor:noarb}; it illustrates the strongest useful meaning of
``universal'' once the benchmark class is specified.

\paragraph{A Hierarchy of Attainable Claims}

The useful alternatives to pathwise universality are not weaker points on one
single order; they change the benchmark, horizon, or probability quantifier.
Table~\ref{tab:hierarchy} records the distinctions.

\begin{table*}[ht]
\caption{\label{tab:hierarchy}Different meanings of a successful or universal strategy.}
\centering
\begin{ruledtabular}
\begin{tabular}{p{0.22\textwidth}p{0.36\textwidth}p{0.34\textwidth}}
\textbf{Notion} & \textbf{Guarantee} & \textbf{Status} \\
\colrule
Pathwise universal profit
& Positive discounted zero-cost gain on every admissible path
& Excluded by classical NA \\
Relative arbitrage
& Outperformance of a specified market or numeraire portfolio
& Possible under structural market assumptions \\
Cover universality
& Vanishing per-period log regret to the best constant-rebalanced portfolio
& Pathwise but benchmark-relative \\
FAPP or statistical edge
& Positive performance under a specified physical law, data protocol, and horizon
& Conditional and empirically falsifiable \\
\end{tabular}
\end{ruledtabular}
\end{table*}

Superhedging duality gives another quantitative notion: once a contingent
claim and an admissible model class are fixed, its superhedging price is the
least initial capital needed for a pathwise guarantee. Requiring zero initial
capital while demanding a strictly positive payoff is precisely the
degenerate, arbitrage-producing endpoint ruled out above.

%%------------------------------------------------------------------------------
\section{FAPP Strategies and Their Honest Scope}
\label{sec:fapp}
%%------------------------------------------------------------------------------

While pathwise universal profit is precluded, strategies may be profitable
\emph{for all practical purposes} (FAPP, following Bell~\cite{bell-a})
relative to a physical measure \(\mathbb{P}\), a benchmark, and a finite
horizon. Market making may earn spreads when adverse selection and inventory
risk are controlled, and Kelly investing maximizes expected logarithmic
growth when the relevant distribution is sufficiently well
specified~\cite{Kelly1956}. Neither is an unconditional profit guarantee.

This is where ``insight'' acquires an operational meaning. A trader specifies
a restricted market class and a mechanism expected to persist in it, then
accepts that the claim is falsifiable outside that class. A lucky finite gain
needs no insight; a repeatable edge, positive conditional expectation, or
sustained benchmark outperformance does.

Economic equilibrium also leaves room for conditional edges. Grossman and
Stiglitz show that perfectly informationally efficient markets are
inconsistent with costly information acquisition: if prices revealed all
information without reward, agents would not pay to produce that
information~\cite{GrossmanStiglitz1980}. Conversely, the Milgrom--Stokey
no-trade theorem shows that, under common priors, common knowledge of
rationality, risk aversion, and an initially Pareto-optimal allocation, private
information alone cannot induce mutually agreeable non-null
trade~\cite{MilgromStokey1982}. These are equilibrium boundaries, not
substitutes for the FTAP.

\paragraph{When the Trader Changes the Market}

The separation between trader and market becomes endogenous when a strategy is
large, crowded, or public. Its orders consume liquidity and move quotes; its
signals invite imitation or anticipatory trading; and its success changes the
distribution on which it was estimated. In that case the deployed price path
is better written \(P^\sigma\): it is partly generated by the use of
\(\sigma\) itself. The analogy is an invasive classical measurement or
feedback experiment, in which the probe perturbs the observed process. In
finance the mechanisms are price impact and strategic adaptation, not a
physical measurement postulate.

Lo's Adaptive Markets Hypothesis~\cite{Lo2004} supplies a less idealized
interpretation. As capital enters a successful strategy, competition, price
impact, and strategic response may erode its edge. The market need not become
the perfect adversary of Theorem~\ref{thm:counterpath}; adaptation merely makes
the physical distribution endogenous to widespread strategy use. Nor does
impact prove that every large strategy loses; it invalidates naive
extrapolation from a backtest that assumes unlimited scale without changing
the environment. Market selection also needs hypotheses: Sandroni shows
conditions under which agents with inaccurate beliefs lose wealth share,
while emphasizing that the result depends on savings and preference
structure~\cite{Sandroni2000}. Selection is therefore an equilibrium
mechanism, not an unconditional theorem that the most accurate forecaster
always dominates.

Empirical strategy discovery adds a separate multiple-testing problem. Trying
many specifications and reporting the best backtest makes its apparent Sharpe
ratio an order statistic rather than an unbiased performance estimate. The
Deflated Sharpe Ratio corrects for selection, non-normal returns, and the
number of trials~\cite{BaileyLopez2014}. Such corrections do not prove an edge;
they make a finite-sample FAPP claim more honestly falsifiable.

Automated execution can amplify feedback when many systems share triggers,
funding constraints, or risk models. The 1987 portfolio-insurance
cascade~\cite{Brady1988} and the 2010 Flash Crash~\cite{Kirilenko2017} illustrate
such interactions. Their relevance comes from market impact and coupled
control rules, not from Rice's theorem.

%%------------------------------------------------------------------------------
\section{Conclusion}
\label{sec:conclusion}
%%------------------------------------------------------------------------------

Everything returns to the same pair: a trader and a market. Fix a total
deterministic trading program and a responsive market can compute the trader's
announced position and reverse its intended effect. This is the diagonal
counterpath and the precise quantifier reversal
\(\forall\sigma\,\exists\mathcal E_\sigma\): every such trader has a market on
which it fails to win. If the trader stays flat, the result is zero; if it
takes a nonzero position, the constructed market makes that trade lose.

Universal computation adds an independent rule-inference boundary. An
arbitrary computation can be encoded either in the behavior of a trader on a
fixed rising path or in a future event of a total computable market generator.
A master program deciding eventual profit or every encoded market event would
therefore decide the Halting Problem. Gold-style inductive inference places
this explicit reduction in a broader learning setting. The claim is not that
restricted market models cannot be learned, but that no computable learner or
certifier covers the unrestricted computationally universal class.

At the passive extreme, a Martin-L\"of-random price-direction record relative
to a computable fair measure need not react to the trader at all. It passes
every effective null test and blocks unbounded success by every
lower-semicomputable nonnegative fair-game capital process. This is not the
adaptive diagonal construction, and neither statement is deterministic chaos.
Between those extremes, a sufficiently large, crowded, or public trader
changes prices and other traders' responses; the trader becomes part of the
market it is trying to exploit.

The existing financial and combinatorial results sharpen, rather than bury,
this core idea. Strict zero-cost discounted profit on every path is classical
arbitrage. Uniform binary averaging fixes expected directional accuracy at
\(1/2\), while optional stopping shows why a high trade win rate need not
imply positive expected gain. The Wheel's put-like payoff, Cover's
benchmark-relative guarantee, equilibrium information results, and FAPP
strategies show what remains possible: conditional growth, relative
outperformance, and transient informational or risk-bearing rents.

A trader can therefore win repeatably only conditionally, by possessing some
model, information, or structural insight that narrows the market class.
Time reversal is the simplest computable falsification probe of that insight;
many other stress transformations are possible. Passing them supports a
limited robustness claim, never an insight-free strategy that wins against
every market. Any practical claim must state its benchmark, financing,
admissibility constraints, horizon, data-selection protocol, and assumptions
about how deployment changes the environment.

\begin{acknowledgments}
Noson S. Yanofsky drew my attention to parallels between the present result,
Arrow's Impossibility Theorem, and fixed-point formulations of
self-reference.

This research was funded in whole or in part by the Austrian Science Fund
(FWF) [Grant DOI: 10.55776/PIN5424624]. For open access purposes, the author
has applied a CC BY public copyright license to any author accepted manuscript
version arising from this submission. The author acknowledges TU Wien
Bibliothek for financial support through its Open Access Funding Programme.

OpenAI Codex (GPT-5) was used to assist with literature organization, LaTeX
editing, and checks of derivations. The author specified the scientific
direction and prompts; model output retained in the manuscript was checked
against the cited sources and explicit calculations. The author accepts full
responsibility for the final text.
\end{acknowledgments}

\section*{Data Availability}
This is a purely mathematical work, and no data or software were created or
analyzed in this study.

%%------------------------------------------------------------------------------
\bibliography{svozil}
\end{document}
